\exercisesection{}

\begin{enumerate}
\def\labelenumi{\arabic{enumi})}
\item
  函数 \(\varphi_0\) 在 \(\mathbf{R}\) 上由 \(\varphi_0(x) = 0\) （若 \(x \leqslant 0\) ）与 \(\varphi_0(x) = e^{-1/x}\) （若 \(x > 0\) ）定义，它是 \(\mathrm{C}^\infty\) 类的，并且满足 \(\varphi_0^{(k)}(0) = 0\) 对每个整数 \(k \in \mathbf{N}\) 。
\item
  设 \(x_0 \in \mathbf{R}^n\) ，\(r > 0\) 且 \(t > 1\) 。存在 \(\varphi \in \mathcal{D}(\mathbf{R}^n)\) 使得
\end{enumerate}

\begin{enumerate}
\def\labelenumi{(\roman{enumi})}
\item
  我们有 \(0 \leqslant \varphi \leqslant 1\) ；
\item
  \(\varphi(x)=1\) 对 \(\|x-x_{0}\|\leqslant r\) 成立，且 \(\varphi(x)=0\) 对 \(\|x-x_{0}\|\geqslant tr\) 成立；
\item
  对每个 \(\alpha \in \mathbf{N}^n\) ，我们有
\end{enumerate}

\[
\sup _ {x \in \mathbf {R} ^ {n}} | \partial^ {\alpha} \varphi (x) | \leqslant \alpha ! ((t - 1) r) ^ {| \alpha |}.
\]

\begin{enumerate}
\def\labelenumi{\arabic{enumi})}
\setcounter{enumi}{2}
\tightlist
\item
  设 \(n \in \mathbf{N}\) 。我们用 \(\mu_n\) 表示 \(\mathbf{R}^n\) 上的勒贝格测度。
\end{enumerate}

\begin{enumerate}
\def\labelenumi{\alph{enumi})}
\tightlist
\item
  设 \(\mathcal{H}_1\) 是映射 \(\partial_1\colon \mathcal{D}(\mathbf{R}^n)\to \mathcal{D}(\mathbf{R}^n)\) 的像；它是满足如下条件的 \(\varphi \in \mathcal{D}(\mathbf{R}^n)\) 所成的空间：
\end{enumerate}

\[
\int_ {\mathbf {R}} \varphi (x, y) d \mu_ {1} (x) = 0
\]

对所有 \(y \in \mathbf{R}^{n-1}\) 成立。

\begin{enumerate}
\def\labelenumi{\alph{enumi})}
\setcounter{enumi}{1}
\tightlist
\item
  设 \(\varphi_0\in \mathcal{D}(\mathbf{R})\) 使其在 \(\mathbf{R}\) 上的积分等于 1。对每个 \(\varphi \in \mathcal{D}(\mathbf{R}^n)\) ，存在 \(\varphi_{1}\in \mathcal{H}_{1}\) 使得
\end{enumerate}

\[
\varphi (x, y) = \varphi_ {1} (x, y) + \varphi_ {0} (x) \int_ {\mathbf {R}} \varphi (t, y) d \mu_ {1} (t)
\]

对所有 \((x,y)\in \mathbf{R}\times \mathbf{R}^{n - 1}\) 成立。

\begin{enumerate}
\def\labelenumi{\alph{enumi})}
\setcounter{enumi}{2}
\tightlist
\item
  对每个广义函数 \(f \in \mathcal{D}'(\mathbf{R}^{n})\) ，存在 \(g \in \mathcal{D}'(\mathbf{R}^{n})\) 使得 \(\partial_{1}g = f\) 。若 \(g_{1}\) 与 \(g_{2}\) 属于 \(\mathcal{D}'(\mathbf{R}^{n})\) 且满足 \(\partial_{1}g_{1} = \partial_{1}g_{2} = f\) ，则存在广义函数 \(h \in \mathcal{D}'(\mathbf{R}^{n-1})\) 使得
\end{enumerate}

\[
\langle g _ {1} - g _ {2}, \varphi \rangle = \langle h, \psi \rangle
\]

对所有 \(\varphi\in\mathcal{D}(\mathbf{R}^{n})\) 成立，其中 \(\psi\in\mathcal{D}(\mathbf{R}^{n-1})\) 由下式定义：

\[
\psi (y) = \int_ {\mathbf {R}} \varphi (t, y) d \mu_ {1} (t).
\]

\(d)\) 设 U 是 \(\mathbf{R}^{n}\) 的连通开子集。设 \(f \in \mathcal{D}'(\mathrm{U})\) 是满足 \(\partial_{i}f = 0\) （对 \(1 \leqslant i \leqslant n\) ）的广义函数。广义函数 \(f\) 是与一个常值函数相伴的广义函数。

\begin{enumerate}
\def\labelenumi{\arabic{enumi})}
\setcounter{enumi}{3}
\tightlist
\item
  我们用 \(\mu\) 表示 \(\mathbf{R}\) 上的勒贝格测度。
\end{enumerate}

\begin{enumerate}
\def\labelenumi{\alph{enumi})}
\tightlist
\item
  设 \(\nu\) 是 \(\mathbf{R}\) 上的测度，\(a \in \mathbf{R}\) 且 \(\nu(\{a\}) = 0\) 。对 \(x \in \mathbf{R}\) 我们令
\end{enumerate}

\[
f (x) = \left\{ \begin{array}{l l} \nu ([ a, x ]) & \text { if } x \geqslant a, \\ \nu ([ x, a ]) & \text { if } x \leqslant a. \end{array} \right.
\]

函数 f 是连续的，并且与之相伴的广义函数满足 \(f' = \nu\) 。

\begin{enumerate}
\def\labelenumi{\alph{enumi})}
\setcounter{enumi}{1}
\item
  设 f 是 R 上的广义函数，使得 \(f'\) 是一个测度 \(\nu\) 。广义函数 f 与 R 上按上一问题方式定义的连续函数相伴。
\item
  设 \(f\) 是 \(\mathbf{R}\) 上的广义函数，使得 \(f^{(k)}\) 对每个整数 \(k \geqslant 0\) 都是测度。那么 \(f\) 是与 \(\mathcal{C}^{\infty}(\mathbf{R})\) 中的函数相伴的广义函数。
\end{enumerate}

\begin{enumerate}
\def\labelenumi{\arabic{enumi})}
\setcounter{enumi}{4}
\item
  与点测度 \(\varepsilon_{n}\) （\(n\in\mathbf{Z}\) ）相伴的广义函数族在 \(\mathcal{S}^{\prime}(\mathbf{R})\) 中可和。其和 \(\omega\) 满足 \(\mathcal{F}(\omega)=\omega\) 。（该公式等价于关于赋有勒贝格测度的局部紧群 \(\mathrm{G}=\mathbf{R}^{n}\) 与子群 \(\mathrm{H}=\mathbf{Z}^{n}\) 的泊松公式，参见 II, § 1, 第 7 目。）
\item
  设 A 是 C 上的交换含单位元代数，满足下列条件：
\end{enumerate}

\begin{enumerate}
\def\labelenumi{(\roman{enumi})}
\item
  代数 \(\mathcal{C}(\mathbf{R})\) 是 A 的含单位元子代数；
\item
  存在导子 \(d\colon\mathrm{A}\to\mathrm{A}\) ，它到 \(\mathcal{C}^{1}(\mathbf{R})\) 的限制与函数的通常导子一致。
\end{enumerate}

\begin{enumerate}
\def\labelenumi{\alph{enumi})}
\item
  R 的恒等映射 X 在 A 中可逆。（证明映射 \(g: \mathbf{R} \to \mathbf{R}\) （它满足 \(g(0) = 0\) ，并且 \(g(x) = x(\log(|x|) - 1)\) 对 \(x \neq 0\) ）是连续的且满足 \(gX = 1\) 。）
\item
  我们有 \(d \circ d(|\mathrm{X}|) = 0\) 。
\item
  不存在双线性映射 \(m\colon\mathcal{D}^{\prime}(\mathbf{R})\times\mathcal{D}^{\prime}(\mathbf{R})\to\mathcal{D}^{\prime}(\mathbf{R})\) 使得 \(m(f,g)=fg\) 对每个偶 \((f,g)\in\mathcal{D}(\mathbf{R})^{2}\) 成立，并且使得
\end{enumerate}

\[
m (f, g) ^ {\prime} = m \left(f ^ {\prime}, g\right) + m \left(f, g ^ {\prime}\right)
\]

对所有 \((f,g)\in\mathcal{D}^{\prime}(\mathbf{R})^{2}\) 成立（“广义函数乘法的不可能性”）。

\begin{enumerate}
\def\labelenumi{\alph{enumi})}
\setcounter{enumi}{3}
\tightlist
\item
  给出 R 上试验函数的序列 \((f_{n})\) 与 \((g_{n})\) 的例子，使得 \(f_{n} \to \varepsilon_{0}\) 且 \(g_{n} \to \varepsilon_{0}\) 在 \(\mathcal{D}'(\mathbf{R})\) 中，并且 \(f_{n}^{2}\) 与 \(g_{n}^{2}\) 在 \(\mathcal{D}'(\mathbf{R})\) 中收敛于不同的极限。
\end{enumerate}

\begin{enumerate}
\def\labelenumi{\arabic{enumi})}
\setcounter{enumi}{6}
\tightlist
\item
  设 \(n \in N\) 。我们用 \(\widetilde{\gamma}\) 表示 \(R^{n}\) 在从 \(R^{n}\) 到 C 的函数所成空间上的线性表示，它由下式定义：
\end{enumerate}

\[
\widetilde {\gamma} (h) f (x) = f (x + h)
\]

对所有 \(h \in \mathbf{R}^n\) 以及每个函数 \(f: \mathbf{R}^n \to \mathbf{C}\) 成立。

\begin{enumerate}
\def\labelenumi{\alph{enumi})}
\item
  对每个 \(h \in R^{n}\) ，映射 \(\widetilde{\gamma}(h)\) 通过向子空间的过渡定义了局部凸空间 \(\mathcal{D}(\mathbf{R}^{n})\) 的一个自同构，它仍记作 \(\widetilde{\gamma}(h)\) 。
\item
  我们用 \(\pmb{\gamma}\) 表示 \(\mathcal{D}'(\mathbf{R}^n)\) 的与 \(\widetilde{\pmb{\gamma}}(-h)\) 转置相应的自同构。若 \(f \in \mathcal{D}(\mathbf{R}^n)\) ，则 \(\pmb{\gamma}(h)f = \widetilde{\pmb{\gamma}}(h)f\) 对所有 \(h \in \mathbf{R}^n\) 成立。
\item
  映射 \(h \mapsto \gamma(h)\) 是 \(\mathbf{R}^{n}\) 在 \(\mathcal{D}'(\mathbf{R}^{n})\) 中的线性表示。
\end{enumerate}

\(d)\) 对每个 \(h \in \mathbf{R}^n\) 以及每个 \(\alpha \in \mathbf{N}^n\) ，我们有 \(\partial^\alpha \circ \boldsymbol{\gamma}(h) = \boldsymbol{\gamma}(h) \circ \partial^\alpha\) 。

\begin{enumerate}
\def\labelenumi{\alph{enumi})}
\setcounter{enumi}{4}
\tightlist
\item
  设 i 是满足 \(1 \leqslant i \leqslant n\) 的整数，\(e_{i}\) 是第 \(i\) 个向量，属于 \(R^{n}\) 的典范基。对每个广义函数 \(f \in \mathcal{D}'(\mathbf{R}^{n})\) ，我们有
\end{enumerate}

\[
\lim_{\substack{h\to 0\\ h\neq 0}}\frac{1}{h} (\boldsymbol {\gamma}(he_{i})f - f) = \partial_{i}f
\]

在赋有弱拓扑的 \(\mathcal{D}'(\mathbf{R}^n)\) 中成立。

\begin{enumerate}
\def\labelenumi{\arabic{enumi})}
\setcounter{enumi}{7}
\item
  设 \(n \in \mathbf{N}\) ，U 是 \(R^{n}\) 的开子集。广义函数 \(f \in \mathcal{D}'(U)\) 称为正的，如果 \(\langle f, \varphi \rangle \geqslant 0\) 对每个 \(\varphi \in \mathcal{D}(U)\) （它满足 \(\varphi \geqslant 0\) ）成立。证明 f 是正的当且仅当 f 是与 U 上一个正测度相伴的广义函数。
\item
  设 \(n \in \mathbf{N}\) ，\(U\) 是 \(\mathbf{R}^n\) 的开子集。设 \(k \in \mathbf{N}\) 且 \(p \in ]1, +\infty[\) 。空间 \(W^{k,p}(U)\) 是自反的。
\item
  设 \(n \in \mathbf{N}\) ，\(U\) 是 \(\mathbf{R}^n\) 的开子集。设 \(k \in \mathbf{N}\) 且 \(p \in [1, +\infty[\)。设 \(V \subset U\) 是相对紧的开子集。
\end{enumerate}

设 \(\varphi \in \mathcal{D}(\mathbf{R}^n)\) 是正函数，其支撑含于 \(\mathbf{R}^n\) 的单位球且积分为 1。对每个 \(\varepsilon > 0\) 以及 \(x \in \mathbf{R}^n\) ，我们令 \(\varphi_{\varepsilon}(x) = \varepsilon^{-n}\varphi(x/\varepsilon)\)。设 \(f \in \mathrm{W}^{k,p}(\mathrm{U})\) 。

\begin{enumerate}
\def\labelenumi{\alph{enumi})}
\item
  \(f_{\mathrm{V}}\) （即 \(f\) 到 V 的限制）属于 \(\mathrm{W}^{k,p}(\mathrm{V})\) 。
\item
  设 \(\widetilde{f}\) 是 \(\mathbf{R}^n\) 上在 U 上与 \(f\) 一致且在 U 外为零的函数。当 \(\varepsilon \to 0\) 时，\(\widetilde{f} * \varphi_{\varepsilon}\) 到 V 的限制属于 \(\mathrm{W}^{k,p}(\mathrm{V})\) 且收敛于 \(f_{\mathrm{V}}\) 在 \(\mathrm{W}^{k,p}(\mathrm{V})\) 中。
\end{enumerate}

\begin{enumerate}
\def\labelenumi{\arabic{enumi})}
\setcounter{enumi}{10}
\tightlist
\item
  设 \(n \in \mathbf{N}\) ，U 是 \(\mathbf{R}^n\) 的开子集。设 \(k \in \mathbf{N}\) 且 \(p \in [1, +\infty[\) 。我们用 \(\mathrm{H}^{k,p}(\mathrm{U})\) 表示 \(\mathcal{C}^{\infty}(\mathrm{U})\) 在 \(\mathrm{W}^{k,p}(\mathrm{U})\) 中的闭包。
\end{enumerate}

对 \(x \in \mathrm{U}\) ，用 \(\delta(x)\) 表示从 \(x\) 到 \(\mathrm{U}\) 在 \(\mathbf{R}^n\) 中的边界的距离；令 \(\mathrm{U}_{-1} = \mathrm{U}_0 = \varnothing\) 并且

\[
\mathrm{U} _ {k} = \{x \in \mathrm{U} | \| x \| <   k, \delta (x) > k ^ {- 1} \}
\]

对每个整数 \(k \geqslant 1\) 。令 \(\mathrm{V}_k = \mathrm{U}_{k+1} \cap (\mathbf{R}^n - \overline{\mathrm{U}}_{k-1})\) 对 \(k \geqslant 1\)。诸开集 \(\mathrm{V}_k\) 覆盖 U；设 \((\psi_k)_{k \geqslant 1}\) 是从属于覆盖 \((\mathrm{V}_k)_{k \geqslant 1}\) 且由 \(\mathcal{D}(\mathrm{U})\) 中的函数组成的单位分解。

\begin{enumerate}
\def\labelenumi{\alph{enumi})}
\tightlist
\item
  对 \(k \geqslant 1\) ，函数
\end{enumerate}

\[
\varphi_{k} = \sum_{\substack{j\geqslant 1\\ \operatorname{Supp}(\psi_{j})\subset \mathrm{V}_{k}}}\psi_{j}
\]

属于 \(\mathcal{D}(\mathrm{U})\) 且支撑含于 \(\mathrm{V}_k\) ；我们有

\[
\sum_ {k \geqslant 1} \varphi_ {k} = 1
\]

在 U 上成立。

\begin{enumerate}
\def\labelenumi{\alph{enumi})}
\setcounter{enumi}{1}
\tightlist
\item
  设 \(f \in \mathrm{W}^{k,p}(\mathrm{U})\) 且 \(\varepsilon > 0\)。设 \(\varphi\) 以及 \(\varphi_{\varepsilon}\) （对 \(\varepsilon > 0\) ）是如习题 10 中那样的函数。对每个 \(k \geqslant 1\)，存在 \(\varepsilon_{k} > 0\) 使得
\end{enumerate}

\[
\left\| \varphi_ {\varepsilon_ {k}} * (\psi_ {k} f) - \psi_ {k} f \right\| _ {k, p} <   \frac {\varepsilon}{2 ^ {k}}.
\]

\begin{enumerate}
\def\labelenumi{\alph{enumi})}
\setcounter{enumi}{2}
\tightlist
\item
  函数
\end{enumerate}

\[
f _ {\varepsilon} = \sum_ {k \geqslant 1} \varphi_ {\varepsilon_ {k}} * (\psi_ {k} f)
\]

属于 \(\mathcal{C}^{\infty}(\mathrm{U})\) 且 \(\| f - f_{\varepsilon}\|_{k,p} < \varepsilon\) 。

\(d)\). 空间 \(\mathrm{H}^{k,p}(\mathrm{U})\) 等于 \(\mathrm{W}^{k,p}(\mathrm{U})\) 。

\begin{enumerate}
\def\labelenumi{\alph{enumi})}
\setcounter{enumi}{4}
\item
  空间 \(\mathrm{H}^{k,p}(\mathrm{U})\) 可等同于 \(\mathcal{C}^{k}(\mathrm{U})\) 关于范数 \(f \mapsto \|f\|_{k,p}\) 的完备化。
\item
  设 V 是 \(R^{n}\) 的开子集，\(\varphi\colon U\to V\) 是 \(C^{k}\) 类微分同胚。假设 \(\leqslant k\) 阶偏导数——\(\varphi\) 的（相应地 \(\varphi^{-1}\) 的）——在 U 上（相应地在 V 上）有界且一致连续。映射 \(f\mapsto f\circ\varphi\) 是从 \(W^{k,p}(V)\) 到 \(W^{k,p}(U)\) 的拓扑向量空间同构。（利用上一问题的结果。）
\end{enumerate}

\begin{enumerate}
\def\labelenumi{\arabic{enumi})}
\setcounter{enumi}{11}
\tightlist
\item
  \begin{enumerate}
  \def\labelenumii{\alph{enumii})}
  \tightlist
  \item
    对所有整数 \(n \in \mathbf{N}\) 与 \(k \in \mathbf{N}\) ，以及每个实数 \(p \geqslant 1\) ，我们有 \(\mathrm{W}^{k,p}(\mathbf{R}^n) = \mathrm{W}_0^{k,p}(\mathbf{R}^n)\) 。
  \end{enumerate}
\end{enumerate}

\begin{enumerate}
\def\labelenumi{\alph{enumi})}
\setcounter{enumi}{1}
\tightlist
\item
  设 \(p \geqslant 1\) 是实数。设 \(U \subset R^{2}\) 是由 \((x, y) \in \mathbf{R}^{2}\) （满足 0 \textless{} \textbar x\textbar{} \textless{} 1 与 0 \textless{} y \textless{} 1）组成的集。特征函数 \(\varphi\) （\(U \cap (\mathbf{R}_{+}^{*} \times \mathbf{R})\) 的）属于 \(W^{1,p}(U)\) ；若 \(\varepsilon > 0\) 充分小，则不存在 \(C^{1}\) 类的函数 f （在 \(\overline{U}\) 的邻域中）使得 \(\|f - \varphi\|_{1,p} < \varepsilon\) 。
\end{enumerate}

\begin{enumerate}
\def\labelenumi{\arabic{enumi})}
\setcounter{enumi}{12}
\tightlist
\item
  设 U 是 \(\mathbf{R}^n\) 的开子集，\(j\colon \mathrm{U}\to \mathbf{R}^n\) 是典范单射。对每个函数 \(f\colon \mathrm{U}\to \mathbf{C}\) ，我们用 \(j_{!}f\) 表示 \(f\) 到 \(\mathbf{R}^n\) 的零延拓，即从 \(\mathbf{R}^n\) 到 \(\mathbf{C}\) 的、在 U 上与 \(f\) 一致且在 U 外为零的函数。
\end{enumerate}

设 \(k \geqslant 0\) 是整数，\(p \geqslant 1\) 是实数。

\begin{enumerate}
\def\labelenumi{\alph{enumi})}
\item
  设 \(f \in \mathrm{W}^{k,p}(\mathrm{U})\) 。对每个 \(\alpha \in \mathbf{N}^n\) （满足 \(|\alpha| \leqslant k\) ），我们有 \(\partial^\alpha(j_{!}f) = j_{!}(\partial^\alpha(f))\) 在 \(\mathcal{D}'(\mathbf{R}^n)\) 中。
\item
  由向子空间的过渡诱导的映射 \(f \mapsto j_{!}f\) 是从 \(\mathrm{W}_{0}^{k,p}(\mathrm{U})\) 到 \(\mathrm{W}^{k,p}(\mathbf{R}^{n})\) 中的等距线性映射。
\end{enumerate}

\begin{enumerate}
\def\labelenumi{\arabic{enumi})}
\setcounter{enumi}{13}
\tightlist
\item
  设 U 是 \(R^{n}\) 的开子集，\(k \geqslant 1\) 是整数。设 \(p \geqslant 1\) 是实数。假设余集 \(R^{n} - U\) 关于勒贝格测度不是可忽略的。用 j 表示 U 到 \(R^{n}\) 中的典范单射。
\end{enumerate}

\begin{enumerate}
\def\labelenumi{\alph{enumi})}
\item
  存在 \(R^{n}\) 中相对紧的开长方体 B，使得 \(B \cap U\) 与 \(B \cap (\mathbf{R}^{n} - \mathbf{U})\) 都不是可忽略的。
\item
  设 \(\widetilde{f} \in \mathcal{D}(\mathbf{R}^n)\) 在 \(B \cap U\) 上等于 1。设 \(f\) 是 \(\widetilde{f}\) 到 U 的限制。那么 \(f\) 不属于 \(\mathrm{W}_0^{k,p}(\mathrm{U})\) （否则，零延拓 \(j_{!}f\) 将属于 \(\mathrm{W}^{k,p}(\mathbf{R}^n)\) ；\(j_{!}f\) 到 B 的限制将是一个满足 \(\partial_i(j_{!}f) = 0\) （对 \(1 \leqslant i \leqslant n\) ）的广义函数，从而将是与一个常数相伴的广义函数。）
\end{enumerate}

\begin{enumerate}
\def\labelenumi{\arabic{enumi})}
\setcounter{enumi}{14}
\tightlist
\item
  设 \(n \geqslant 1\) 是整数，p 是满足 \(1 \leqslant p < n\) 的实数。
\end{enumerate}

\begin{enumerate}
\def\labelenumi{\alph{enumi})}
\item
  设 \(\alpha\) 是满足 \(0 < \alpha < n/p - 1\) 的实数。设 \(g\) 是 \(\mathbf{R}_{+}^{*}\) 上的无穷可微正函数，使得 \(g(x) = x^{-\alpha}\) （当 \(x\) 充分小）且 \(g(x) = 0\) （当 \(x\) 充分大）；从 \(\mathbf{R}^{n}\) 到 \(\mathbf{C}\) 的、由 \(x \mapsto g(\|x\|)\) 定义的函数属于 \(\mathrm{W}^{1,p}(\mathbf{R}^{n})\) 。
\item
  设 \(i \mapsto q_{i}\) 是从 N 到 \(Q^{n}\) 的双射，f 是由下式定义的从 \(R^{n}\) 到 C 的映射：
\end{enumerate}

\[
f (x) = \sum_ {i \in \mathbf {N}} 2 ^ {- i} g (\| x - q _ {i} \|)
\]

对 \(x \in \mathbf{R}^n\) 成立。定义 \(f\) 的级数在 \(\mathrm{W}^{1,p}(\mathbf{R}^n)\) 中收敛；对 \(\mathbf{R}^n\) 的每个开集 U，\(f\) 到 U 的限制都不是几乎处处可微的。\(c)\) 存在在任一点都不连续的函数 \(f \in \mathrm{W}^{1,n}(\mathbf{R}^n)\) 。（对函数 \(x \mapsto \log(|\log(|x|)|)\) 用类似的方法。）

\begin{enumerate}
\def\labelenumi{\arabic{enumi})}
\setcounter{enumi}{15}
\tightlist
\item
  设 \(n \geqslant 1\) 与 \(k \geqslant 1\) 是整数。设 \(p \geqslant 1\) 是实数。设 U 是 \(\mathbf{R}^n\) 的有界开子集。
\end{enumerate}

\begin{enumerate}
\def\labelenumi{\alph{enumi})}
\item
  假设 \(n \neq kp\) 。设 q 是满足 \(1 \leqslant q \leqslant np/(n-kp)\) 的实数。\(\mathcal{D}(U)\) 到 \(L^{q}(U)\) 中的单射容许一个连续延拓 \(i\) （从 \(W_{0}^{k,p}(U)\) 到 \(L^{q}(U)\) 中）。
\item
  假设 \(1 \leqslant q \leqslant np/(n-kp)\) 。线性映射 \(i\) 是紧的（“雷利希–孔德拉绍夫定理”）。
\item
  假设 \(k \geqslant 1\) 且 \(q = np/(n - p)\) 。线性映射 i 不是紧的。
\end{enumerate}

\begin{enumerate}
\def\labelenumi{\arabic{enumi})}
\setcounter{enumi}{16}
\tightlist
\item
  设 \(n \geqslant 1\) 是整数。我们用 \(\mu\) 表示 \(\mathbf{R}^n\) 上的勒贝格测度。
\end{enumerate}

\begin{enumerate}
\def\labelenumi{\alph{enumi})}
\tightlist
\item
  \(\mathrm{L}^{2}(\mathbf{R}^{n},\mu)\) 的有界子集 B 是紧的，当且仅当对每个 \(\varepsilon>0\) ，存在 \(R^{n}\) 的紧子集 K 使得
\end{enumerate}

\[
\int_ {\mathbf {R} ^ {n} - K} | f | ^ {2} + \int_ {\mathbf {R} ^ {n} - K} | \widehat {f} | ^ {2} <   \varepsilon
\]

对所有 \(f\in \mathrm{B}\) 成立。

\begin{enumerate}
\def\labelenumi{\alph{enumi})}
\setcounter{enumi}{1}
\tightlist
\item
  设 \(k \in N\) 。设 K 是 \(R^{n}\) 的紧子集。我们用 \(\mathrm{H}_{k}^{2}(\mathrm{K})\) 表示索博列夫空间 \(\mathrm{H}_{k}^{2}(\mathbf{R}^{n})\) 的由在紧集 K 外为零的函数组成的闭子空间。典范单射 \(\mathrm{H}_{k}^{2}(\mathrm{K}) \to \mathrm{H}_{k}^{2}(\mathbf{R}^{n})\) 是紧的。
\end{enumerate}

\begin{enumerate}
\def\labelenumi{\arabic{enumi})}
\setcounter{enumi}{17}
\tightlist
\item
  设 \(n \in \mathbf{N}\) ，\(U\) 是 \(\mathbf{R}^n\) 的开子集。本习题的目标是证明 \(\mathcal{D}'(U)\) 是有界型空间。
\end{enumerate}

若 K 是 U 的紧子集，我们给 \(\mathcal{C}_{\mathrm{K}}^{\infty}(\mathrm{U})'\) 赋予有界收敛拓扑。我们用 \(\pi_{\mathrm{K}}\colon\mathcal{D}'(\mathrm{U})\to\mathcal{C}_{\mathrm{K}}^{\infty}(\mathrm{U})'\) 表示把广义函数映为其到 \(\mathcal{C}_{\mathrm{K}}^{\infty}(\mathrm{U})\) 的限制的线性映射；它是连续的。

设 \(S \subset \mathcal{D}'(U)\) 是凸的、均衡的且吸收有界集的子集。

\begin{enumerate}
\def\labelenumi{\alph{enumi})}
\item
  存在 U 的紧子集 K 使得 \(\mathrm{Ker}(\pi_{\mathrm{K}}) \subset \mathrm{S}\) 。（否则，考虑 U 的紧子集的递增序列 \((\mathrm{K}_{m})_{m \in \mathbf{N}}\) ，其诸内部覆盖 U，并选取 \(f_{m}\) 于 \(\mathrm{Ker}(\pi_{\mathrm{K}_{m}}) - \mathrm{S}\) 中；那么诸元素 \(mf_{m}\) （\(m \in N\) ）所成之集是一个不被 S 吸收的有界子集。）
\item
  由广义函数 \(f \in \mathcal{D}'(\mathrm{U})\) （满足 \(f + \operatorname{Ker}(\pi_{\mathrm{K}}) \subset \mathrm{S}\) ）组成的集 T 是凸的、均衡的且吸收有界集的集。
\item
  集 \(\pi_{\mathrm{K}}(\mathrm{T})\) 是 \(\mathcal{C}_{\mathrm{K}}^{\infty}(\mathrm{U})'\) 的凸的、均衡的且吸收有界集的子集。（设 A 是 \(\mathcal{C}_{\mathrm{K}}^{\infty}(\mathrm{U})'\) 的有界子集；证明存在 \(\mathcal{D}(\mathrm{U})\) 中 0 的邻域 W 使得 \(\mathrm{A} \subset (\mathrm{W} \cap \mathcal{C}_{\mathrm{K}}^{\infty}(\mathrm{U}))^{\circ}\) ，并由此推出 A 含于 \(\pi_{\mathrm{K}}(\mathrm{W}^{\circ})\) ，从而被 \(\pi_{\mathrm{K}}(\mathrm{T})\) 吸收。）
\end{enumerate}

\(d)\) 由此推出 S 是 0 的邻域，并断言空间 \(\mathcal{D}'(U)\) 是有界型的。（注意 \(\pi_{\mathrm{K}}(\mathrm{T})\) 是 0 的邻域，并注意 S 包含 \(\pi_{\mathrm{K}}^{-1}(\pi_{\mathrm{K}}(\mathrm{T}))\) 。）

\begin{enumerate}
\def\labelenumi{\arabic{enumi})}
\setcounter{enumi}{18}
\tightlist
\item
  设 \(n \in \mathbf{N}\) ，\(\mathrm{U} \subset \mathbf{R}^n\) 是开集。
\end{enumerate}

\begin{enumerate}
\def\labelenumi{\alph{enumi})}
\tightlist
\item
  设 \(f \in \mathcal{D}'(\mathrm{U})\) 。开集 \(\mathrm{W} \subset \mathrm{U}\) （使得 \(f\) 到 W 的限制为零）的并 V 在 U 中是开的。
\end{enumerate}

f 的支撑称为 V 的余集；它是 U 的闭子集。

\begin{enumerate}
\def\labelenumi{\alph{enumi})}
\setcounter{enumi}{1}
\item
  设 \(f \in \mathcal{D}'(\mathrm{U})\) ，\(\varphi \in \mathcal{D}(\mathrm{U})\) 的支撑与 f 的支撑不相交。我们有 \(\langle f, \varphi \rangle = 0\) 。
\item
  设 \(f \in \mathcal{D}'(\mathrm{U})\) 且 \(\alpha \in N^{n}\) 。\(\partial^{\alpha}f\) 的支撑含于 f 的支撑。
\item
  若 \(f \in \mathcal{D}'(\mathrm{U})\) 是与测度 \(\nu\) 相伴的广义函数，则 f 的支撑与 \(\nu\) 的支撑重合。
\item
  设 \(f\) 是 U 中具有紧支撑的广义函数。线性形式 \(f\) 可延拓为空间 \(\mathcal{C}^{\infty}(\mathrm{U})\) 上的连续线性形式。
\item
  U 中具有紧支撑的广义函数所成的空间可等同于空间 \(\mathcal{C}^{\infty}(\mathrm{U})\) 的对偶。
\end{enumerate}

\(g)\) 设 \(f \in \mathcal{D}'(\mathbf{R}^n)\) 是具有紧支撑的广义函数。我们有 \(f \in \mathcal{S}'(\mathbf{R}^n)\) 。\(h)\) 设 \(f \in \mathcal{D}'(\mathbf{R}^n)\) 是具有紧支撑的广义函数。\(f\) 的傅里叶变换属于 \(\mathcal{C}^\infty(\mathbf{R}^n)\) ；它是连同其所有导数都具有多项式增长的函数。

\begin{enumerate}
\def\labelenumi{\arabic{enumi})}
\setcounter{enumi}{19}
\tightlist
\item
  设 \(n \in \mathbf{N}\) 。我们用 \(V_n\) 表示 \(\mathbf{C}\) 的由复数 \(\lambda\) （它们满足 \(\mathcal{R}(\lambda) > -n - 1\) 且 \(\lambda \notin \{-1, -2, \ldots, -n\}\) ）组成的开集。
\end{enumerate}

设 \(\lambda\in V_{0}\) 。我们用 \(x_{+}^{\lambda}\) 表示 \(\mathbf{R}\) 上的由局部可积函数

\[
x \mapsto \left\{ \begin{array}{l l} 0 & \text { if } x \leqslant 0, \\ x ^ {\lambda} & \text { if } x > 0. \end{array} \right.
\]

\begin{enumerate}
\def\labelenumi{\alph{enumi})}
\item
  映射 \(\lambda \mapsto x_{+}^{\lambda}\) （从 \(\mathrm{V}_0\) 到 \(\mathcal{D}'(\mathbf{R})\) 中的）是解析的。
\item
  设 \(n \in \mathbf{N}\) 。对每个 \(\lambda \in V_n\) ，在 \(\mathcal{D}(\mathbf{R})\) 上由
\end{enumerate}

\[
\begin{array}{l} \varphi \mapsto \int_ {0} ^ {1} x ^ {\lambda} \Big (\varphi (x) - \varphi (0) - x \varphi^ {\prime} (0) - \dots - \frac {x ^ {n - 1}}{(n - 1) !} \varphi^ {(n - 1)} (0) \Big) d x \\ \qquad + \int_ {1} ^ {+ \infty} x ^ {\lambda} \varphi (x) d x + \sum_ {k = 1} ^ {n} \frac {\varphi^ {(k - 1)} (0)}{(k - 1) ! (\lambda + k)} \end{array}
\]

定义的映射是缓增广义函数；它不依赖于使 \(\lambda \in V_{n}\) 的整数 n，并且与 \(x_{+}^{\lambda}\) 一致当 \(\lambda \in V_{0}\) 时。

我们将用 \(x_{+}^{\lambda}\) 表示这样定义的广义函数，对每个 \(\lambda\) 属于诸集 \(V_{n}\) （\(n \in N\) ）之并。

\begin{enumerate}
\def\labelenumi{\alph{enumi})}
\setcounter{enumi}{2}
\tightlist
\item
  设 \(n \in N\) ，\(\lambda \in C\) 满足 \(-n - 1 < \mathcal{R}(\lambda) < -n\)。对 R 上每个施瓦茨函数 \(\varphi\) ，我们有
\end{enumerate}

\[
\langle x _ {+} ^ {\lambda}, \varphi \rangle = \int_ {0} ^ {+ \infty} x ^ {\lambda} \left(\varphi (x) - \varphi (0) - x \varphi^ {\prime} (0) - \dots - \frac {x ^ {n - 1}}{(n - 1) !} \varphi^ {(n - 1)} (0)\right) d x.
\]

  从 \(V_{n}\) 到 \(\mathcal{D}'(\mathrm{U})\) 中的由 \(\lambda\mapsto x_{+}^{\lambda}\) 定义的映射是解析的。

\(d)\) 对每个 \(\lambda\) （属于诸集 \(V_{m}\) 之并的），我们有公式

\[
(x _ {+} ^ {\lambda}) ^ {\prime} = \lambda x _ {+} ^ {\lambda - 1}.
\]

\begin{enumerate}
\def\labelenumi{\arabic{enumi})}
\setcounter{enumi}{20}
\tightlist
\item
  \begin{enumerate}
  \def\labelenumii{\alph{enumii})}
  \tightlist
  \item
    设 \(f \in \mathcal{S}(\mathbf{R})\) 。R 上满足 \(g(0) = 2f'(0)\) 且 \(g(x) = (f(x) - f(-x))/x\) （对所有 \(x \neq 0\) ）的函数 g 在 R 上关于勒贝格测度可积。把 f 映为 g 在 R 上的积分的映射是一个缓增广义函数，记作 vp，称为 1/x 的柯西主值。
  \end{enumerate}
\end{enumerate}

\begin{enumerate}
\def\labelenumi{\alph{enumi})}
\setcounter{enumi}{1}
\tightlist
\item
  对每个 \(f \in \mathcal{S}(\mathbf{R})\) ，我们有
\end{enumerate}

\[
\langle \mathrm{vp}, f \rangle = \lim _ {\varepsilon \rightarrow 0} \int_ {\mathbf {R} - [ - \varepsilon , \varepsilon ]} \frac {f (x)}{x} d x.
\]

\begin{enumerate}
\def\labelenumi{\alph{enumi})}
\setcounter{enumi}{2}
\item
  vp 到 \(\mathbf{R} - \{0\}\) 的限制是与局部可积函数 \(x \mapsto 1/x\) （\(\mathbf{R} - \{0\}\) 上的）相伴的广义函数。
\item
  广义函数 vp 与下述广义函数的导数重合：该广义函数与 R 上在 0 处取值 0 且对 \(x \neq 0\) 由 \(x \mapsto \log(|x|)\) 定义的局部可积函数相伴。
\item
  设 \(\varepsilon > 0\) 。我们用 \(p_{\varepsilon}\) 表示函数 \(x \mapsto 1 / (x + i\varepsilon)\) （\(\mathbf{R}\) 上的）。我们有
\end{enumerate}

\[
\lim _ {\varepsilon \to 0} p _ {\varepsilon} = \mathrm{vp} - i \pi \varepsilon_ {0}
\]

在赋有弱拓扑的空间 \(\mathcal{S}'(\mathbf{R})\) 中成立，其中 \(\varepsilon_0\) 是与 0 处的点质量相伴的缓增广义函数。

\begin{enumerate}
\def\labelenumi{\arabic{enumi})}
\setcounter{enumi}{21}
\item
  R 上与指数函数相伴的广义函数不是缓增的。
\item
  设 \(n \in \mathbf{N}\) ，U 是 \(\mathbf{R}^n\) 的开集。我们用 \(\mu\) 表示 U 上的勒贝格测度。对 \(m \in \mathbf{N}\) ，我们记 \(\boldsymbol{m} = (m, \ldots, m) \in \mathbf{N}^n\) 。
\end{enumerate}

\begin{enumerate}
\def\labelenumi{\alph{enumi})}
\tightlist
\item
  设 \(f \in \mathcal{D}(\mathrm{U})\) 。假设存在整数 \(m \geqslant 1\) 以及实数 \(c \geqslant 0\) 使得
\end{enumerate}

\[
| \langle f, \varphi \rangle | \leqslant c \int_ {\mathrm{U}} | \partial^ {\boldsymbol {m}} \varphi | d \mu
\]

对所有 \(\varphi\in\mathcal{D}(\mathrm{U})\) 成立。于是存在函数 \(g\in\mathrm{L}^{\infty}(\mathrm{U})\) 使得 \(f=\partial^{m}g\) 。（线性映射 \(\varphi\mapsto\partial^{m}\varphi\) 在 \(\mathcal{D}(\mathrm{U})\) 上是单射；设 E 是它的像；考虑 E 上把 \(\partial^{m}\varphi\) 映为 \(\langle f,\varphi\rangle\) 的到 C 中的线性形式，并应用哈恩–巴拿赫定理。）

\begin{enumerate}
\def\labelenumi{\alph{enumi})}
\setcounter{enumi}{1}
\tightlist
\item
  设 V 是含于 U 的相对紧开集。存在整数 \(m \geqslant 1\) 以及 \(c \geqslant 0\) 使得
\end{enumerate}

\[
| \langle f, \varphi \rangle | \leqslant c \int_ {\mathrm{U}} | \partial^ {\boldsymbol {m}} \varphi | d \mu
\]

对所有支撑含于 V 的 \(\varphi\in\mathcal{D}(\mathrm{U})\) 成立。

\begin{enumerate}
\def\labelenumi{\alph{enumi})}
\setcounter{enumi}{2}
\tightlist
\item
  对每个相对紧开集 \(V \subset U\) ，存在连续函数 \(g \in \mathcal{C}(V)\) 以及 \(\alpha \in N^{n}\) ，使得 f 到 V 的限制等于 \(\partial^{\alpha}g\)。\(d)\) 给出一个广义函数 \(f \in \mathcal{D}'(\mathbf{R}^{n})\) 的例子，使得 f 不是 \(f = \partial^{\alpha}g\) 的形式，其中 \(g \in \mathcal{C}(\mathbf{R}^{n})\) 是连续函数且 \(\alpha \in N^{n}\) 。
\end{enumerate}

\begin{enumerate}
\def\labelenumi{\arabic{enumi})}
\setcounter{enumi}{23}
\tightlist
\item
  设 \(n \in N\) ，U 是 \(R^{n}\) 的开集。设 \(f \in \mathcal{D}'(U)\) 是支撑 C 在 U 中为紧的广义函数。
\end{enumerate}

\begin{enumerate}
\def\labelenumi{\alph{enumi})}
\tightlist
\item
  存在整数 \(m \in N\) ，使得广义函数 f 延拓为空间 \(\mathcal{C}^{m}(U)\) 上的连续线性形式，该空间赋有函数及其阶 \(\leqslant m\) 的偏导数在紧子集上一致收敛的拓扑。
\end{enumerate}

满足该性质的最小整数 m 称为 f 的阶。

\begin{enumerate}
\def\labelenumi{\alph{enumi})}
\setcounter{enumi}{1}
\item
  设 \(\varphi \in \mathcal{D}(\mathrm{U})\) 使得 \(\partial^{\alpha}\varphi\) 对每个 \(\alpha \in \mathbf{N}^{n}\) （满足 \(|\alpha| \leqslant m\) ）在 C 上为零。我们则有 \(\langle f, \varphi \rangle = 0\) 。（对充分小的 \(\delta > 0\) ，考虑 C 的邻域 \(\mathrm{U}_{\delta}\) ，它由 \(x \in \mathrm{U}\) （满足 \(d(x, \mathrm{C}) \leqslant \delta\) ）组成，并证明存在 \(\eta \geqslant 0\) （当 \(\delta\) 趋于 0 时趋于 0），使得对某个 \(\mathrm{C}_{0} \geqslant 0\) ，有 \(|\partial^{\alpha}\varphi(x)| \leqslant \mathrm{C}_{0}\delta^{m-|\alpha|}\eta\) 对 \(x \in \mathrm{U}_{\delta}\) 与 \(|\alpha| \leqslant m\) ；证明存在函数 \(\beta_{\delta} \in \mathcal{D}(\mathrm{U})\) ，它在 \(\mathrm{U}_{\delta/4}\) 上等于 1 且支撑含于 \(\mathrm{U}_{\delta}\) ，并满足 \(|\partial^{\alpha}\beta_{\delta}| \leqslant \mathrm{C}_{1}\delta^{-|\alpha|}\) 对某个 \(\mathrm{C}_{1} \geqslant 0\) ；最后注意 \(\langle f, \varphi \rangle = \langle f, \varphi\beta_{\delta} \rangle\) 且 \(\varphi\beta_{\delta}\) 在 \(\mathcal{C}^{m}(\mathrm{U})\) 中趋于 0。）
\item
  设 \(x \in U\) 。支撑化为单点 x 的广义函数 f 所成的空间以诸广义函数 \(\partial^{\alpha}\varepsilon_{x}\) （其中 \(\alpha \in N^{n}\) ）为代数基。
\item
  设 \(n = 1\) 。公式
\end{enumerate}

\[
\langle f, \varphi \rangle = \lim _ {m \rightarrow + \infty} \left(\sum_ {1 \leqslant k \leqslant m} \varphi (k ^ {- 1}) - m \varphi (0) - \log (m) \varphi^ {\prime} (0)\right)
\]

在 \(\mathbf{R}\) 上定义了一个阶为 1 的广义函数，其支撑是紧的，等于集 \(C = \{0\} \cup \{1/(n+1) \mid n \in \mathbf{N}\}\) 。存在序列 \((\varphi_n)\) （在 \(\mathcal{D}(\mathbf{R})\) 中），使得 \(\partial^{\alpha}\varphi_{n}\) 对每个 \(\alpha\in\mathbf{N}^{n}\) （满足 \(|\alpha|\leqslant1\) ）在 C 上一致收敛于 0，但 \(\langle f,\varphi_{n}\rangle\to+\infty\) 。

\begin{enumerate}
\def\labelenumi{\arabic{enumi})}
\setcounter{enumi}{24}
\tightlist
\item
  设 \(n \in \mathbf{N}\) 。我们用 \(N\) 表示 \(\mathbf{R}^n\) 上的欧几里得范数，并用 \(\mu\) 表示勒贝格测度。
\end{enumerate}

我们用 \(\mathcal{D}_{\mathrm{L}^{1}}(\mathbf{R}^{n})\) 表示 \(\mathcal{C}^{\infty}(\mathbf{R}^{n})\) 的向量子空间，其元素是函数 \(\varphi\) ，它满足 \(\partial^{\alpha}\varphi\in\mathrm{L}^{1}(\mathbf{R}^{n})\) 且 \(\partial^{\alpha}\varphi\in\mathcal{C}_{0}(\mathbf{R}^{n})\) 对所有 \(\alpha\in\mathbf{N}^{n}\) 。\(a)\) 空间 \(\mathcal{D}_{\mathrm{L}^{1}}(\mathbf{R}^{n})\) 赋有由半范数

\[
q _ {\alpha} (\varphi) = \int_ {\mathbf {R} ^ {n}} | \partial^ {\alpha} \varphi | d \mu , \quad \alpha \in \mathbf {N} ^ {n},
\]

定义的拓扑后是弗雷歇空间。

\begin{enumerate}
\def\labelenumi{\alph{enumi})}
\setcounter{enumi}{1}
\item
  空间 \(\mathcal{D}(\mathbf{R}^n)\) 含于 \(\mathcal{D}_{\mathrm{L}^1}(\mathbf{R}^n)\) 且在其中稠密。可以把对偶空间 \(\mathcal{D}_{\mathrm{L}^1}'(\mathbf{R}^n)\) （\(\mathcal{D}_{\mathrm{L}^1}(\mathbf{R}^n)\) 的）等同于 \(\mathcal{D}'(\mathbf{R}^n)\) 的一个子空间。
\item
  设 \(f \in \mathcal{D}'(\mathbf{R}^n)\) 。我们有 \(f \in \mathcal{D}_{\mathrm{L}^1}'(\mathbf{R}^n)\) 当且仅当存在有限集 I、族 \((\alpha_i)_{i \in \mathrm{I}}\) （在 \(\mathbf{N}^n\) 中）以及族 \((g_i)_{i \in \mathrm{I}}\) （在 \(\mathrm{L}^\infty(\mathbf{R}^n)\) 中）使得
\end{enumerate}

\[
f = \sum_ {i \in I} \partial^ {\alpha_ {i}} g _ {i}.
\]

（确定 I 与族 \((\alpha_{i})\) ，使得 \(|\langle f,\varphi\rangle |\) 有界，当 \(\varphi\) 满足 \(q_{\alpha_i}(\varphi)\leqslant 1\) 对所有 \(i\in \mathrm{I}\) ；把 \(\mathcal{D}_{\mathrm{L}^1}(\mathbf{R}^n)\) 等同于 \(\mathrm{L}^1 (\mathbf{R}^n)^m\) 的一个向量子空间，借助映射 \(\varphi \mapsto (\partial^{\alpha_i}\varphi)_{i\in \mathrm{I}}\) ，并应用哈恩–巴拿赫定理。）

\begin{enumerate}
\def\labelenumi{\alph{enumi})}
\setcounter{enumi}{3}
\tightlist
\item
  设 \(f \in \mathcal{D}'(\mathbf{R}^n)\) 。我们有 \(f \in \mathcal{D}_{\mathrm{L}^1}'(\mathbf{R}^n)\) 当且仅当存在具有缓和增长的连续函数 \(g\) （在 \(\mathbf{R}^n\) 上）以及 \(\alpha \in \mathbf{N}^n\) 使得 \(f = \partial^\alpha g\)。\(e)\) 设 \(f \in \mathcal{D}'(\mathbf{R}^n)\) 。存在整数 \(m \in \mathbf{N}\) 使得 \((1 + \mathrm{N}^2)^{-m/2}f\) 属于 \(\mathcal{D}_{\mathrm{L}^1}'(\mathbf{R}^n)\) 。
\end{enumerate}

\(f)\) 设 \(f \in \mathcal{D}'(\mathbf{R}^n)\) 。我们有 \(f \in \mathcal{S}'(\mathbf{R}^n)\) 当且仅当存在有界连续函数 \(g\colon \mathbf{R}^n \to \mathbf{C}\) 、整数 \(m \in \mathbf{N}\) 以及 \(\alpha \in \mathbf{N}^n\) 使得 \(f = \partial^\alpha((1 + \mathrm{N}^2)^{m/2}g)\) 。

\begin{enumerate}
\def\labelenumi{\arabic{enumi})}
\setcounter{enumi}{25}
\tightlist
\item
  设 \(n \in N\) 。我们给 \(R^{n}\) 赋予勒贝格测度 \(\mu\) ，并把 \(R^{n}\) 与其对偶通过映射 \((x, y) \mapsto \exp(2i\pi x \cdot y)\) 等同。
\end{enumerate}

\begin{enumerate}
\def\labelenumi{\alph{enumi})}
\item
  对任何函数 \(f \in \mathcal{S}(\mathbf{R}^n)\) ，映射 \(g \mapsto f * g\) 是局部凸空间 \(\mathcal{S}(\mathbf{R}^n)\) 的自同构。
\item
  设 \(f \in \mathcal{S}'(\mathbf{R}^{n})\) 且 \(g \in \mathcal{S}(\mathbf{R}^{n})\) 。映射 \(\varphi \mapsto \langle f, \check{g} * \varphi \rangle\) 是一个缓增广义函数，称为 f 与 g 的卷积。若 \(f \in \mathcal{S}(\mathbf{R}^{n})\) ，它与同 \(f * g\) 相伴的广义函数一致。我们仍用 \(f * g\) 表示这样定义的缓增广义函数。
\end{enumerate}

在下面的问题中，我们固定 \(f \in \mathcal{S}'(\mathbf{R}^n)\) 与 \(g \in \mathcal{S}(\mathbf{R}^n)\) 。\(c)\) 我们有 \(f * g \in \mathcal{C}^\infty(\mathbf{R})\) 并且

\[
\partial^ {\alpha} (f * g) = \partial^ {\alpha} (f) * g = f * \partial^ {\alpha} (g)
\]

对所有 \(\alpha\in\mathbf{N}^{n}\) 成立。

\(d)\) 对每个 \(y \in \mathbf{R}^n\)，用 \(g_y\) 表示函数 \(x \mapsto g(y - x)\) （\(\mathbf{R}^n\) 上的）。函数 \(f * g\) 满足

\[
(f * g) (y) = \langle f, g _ {y} \rangle
\]

对所有 \(y \in R^{n}\) 成立。

\begin{enumerate}
\def\labelenumi{\alph{enumi})}
\setcounter{enumi}{4}
\item
  设 \(h \in \mathcal{S}(\mathbf{R}^n)\) 。我们有 \((f*g)*h = f*(g*h)\) 。
\item
  函数 \(f*g\) 具有多项式增长，其所有导数亦然。（利用习题 25。）
\item
  我们有 \(\mathcal{F}(f*g) = \mathcal{F}(f)\mathcal{F}(g)\) 。
\item
  fg 的傅里叶变换是 \(R^{n}\) 上的无穷可微函数，并且连同其所有导数都具有多项式增长。此外，对每个 \(y \in R^{n}\) ，我们有
\end{enumerate}

\[
\mathcal {F} (f g) (y) = \langle f, g e _ {- y} \rangle
\]

其中 \(e_y\) 表示特征标 \(x \mapsto \exp(2i\pi x \cdot y)\) （\(\mathbf{R}^n\) 的）。

\begin{enumerate}
\def\labelenumi{\arabic{enumi})}
\setcounter{enumi}{26}
\tightlist
\item
  设 \(n \in \mathbf{N}\)。我们用 N 表示 \(\mathbf{R}^{n}\) 上的欧几里得范数，并用 \(\mu\) 表示 \(\mathbf{R}^{n}\) 的每个开子集上的勒贝格测度。设 \(\Delta\) 是微分算子
\end{enumerate}

\[
\Delta = - \sum_ {i = 1} ^ {n} \partial_ {i} ^ {2}
\]

将它视为 \(\mathcal{D}'(\mathbf{R}^{n})\) 的自同态。

对 \(k \in \mathbf{R}\) ，我们用 \(\mathrm{S}^k(\mathbf{R}^n)\) 表示这样的缓增广义函数 \(f\) （\(\mathcal{S}'(\mathbf{R}^n)\) 中的）所成的空间：\((1 + \mathrm{N}^k)\mathcal{F}(f)\) 属于 \(\mathrm{L}^2(\mathbf{R}^n)\) 。

设 \(U \subset R^{n}\) 是 \(R^{n}\) 的开子集。我们用 \(\mathrm{S}_{\mathrm{U}}^{k}(\mathbf{R}^{n})\) 表示满足如下条件的广义函数 \(f \in \mathcal{D}'(\mathbf{R}^{n})\) 所成的空间：对每个支撑含于 U 的试验函数 \(\varphi \in \mathcal{D}(\mathbf{R}^{n})\) ，我们有 \(\varphi f \in \mathrm{S}^{k}(\mathbf{R}^{n})\) 。

\begin{enumerate}
\def\labelenumi{\alph{enumi})}
\tightlist
\item
  空间 \(\mathrm{S}^{k}(\mathbf{R}^{n})\) 赋有内积
\end{enumerate}

\[
\langle f _ {1} \mid f _ {2} \rangle = \int_ {\mathbf {R} ^ {n}} (1 + \mathrm{N} ^ {k}) \overline {{\mathcal {F} (f _ {1})}} \mathcal {F} (f _ {2}) d \mu .
\]

当 \(k \in \mathbf{N}\) 时，我们有 \(W^{k}(\mathbf{R}^{n}) = S^{k}(\mathbf{R}^{n})\) ，并且这两个巴拿赫空间的范数等价；此外，这时我们有 \(W^{k}(\mathbf{R}^{n}) \subset S_{\mathrm{U}}^{k}(\mathbf{R}^{n})\) 。

\begin{enumerate}
\def\labelenumi{\alph{enumi})}
\setcounter{enumi}{1}
\tightlist
\item
  \(\mathrm{S}^{k}(\mathbf{R}^{n})\) 的对偶通过双线性形式等同于 \(\mathrm{S}^{-k}(\mathbf{R}^{n})\)
\end{enumerate}

\[
(f, g) \mapsto \int_ {\mathbf {R} ^ {n}} \mathcal {F} (f) (- x) \mathcal {F} (g) (x) d \mu (x)
\]

（对 \((f,g)\in \mathrm{S}^{-k}(\mathbf{R}^n)\times \mathrm{S}^k (\mathbf{R}^n).\)）

\begin{enumerate}
\def\labelenumi{\alph{enumi})}
\setcounter{enumi}{2}
\tightlist
\item
  假设 U 是相对紧的。我们有
\end{enumerate}

\[
\bigcup_ {k \in \mathbf {R}} \mathrm{S} _ {\mathrm{U}} ^ {k} (\mathbf {R} ^ {n}) = \mathcal {D} ^ {\prime} (\mathbf {R} ^ {n}).
\]

（设 \(f \in \mathcal{D}'(\mathbf{R}^n)\) ；把习题 19 的结果应用于 \(\varphi f\) ，其中 \(\varphi \in \mathcal{D}(\mathbf{R}^n)\) 在 U 上等于 1。）

\(d)\) 设 \(k \in \mathbf{R}\) 且 \(f \in \mathrm{S}^k(\mathbf{R}^n)\) 。对每个整数 \(i\) （满足 \(1 \leqslant i \leqslant n\) ），我们有 \(\partial_i f \in \mathrm{S}^{k-1}(\mathbf{R}^n)\) 。此外，若 \(\Delta(f) \in \mathrm{S}^k(\mathbf{R}^n)\) ，则 \(f \in \mathrm{S}^{k+2}(\mathbf{R}^n)\) 。

\begin{enumerate}
\def\labelenumi{\alph{enumi})}
\setcounter{enumi}{4}
\tightlist
\item
  设 \(k \in \mathbf{R}\) 且 \(f \in \mathrm{S}_{\mathrm{U}}^{k}(\mathbf{R}^{n})\) 。若 \(\Delta(f) \in \mathrm{S}_{\mathrm{U}}^{k}(\mathbf{R}^{n})\) ，则 \(f \in \mathrm{S}_{\mathrm{U}}^{k+2}(\mathbf{R}^{n})\) 。
\end{enumerate}

\(f)\) 设 \(k > n/2\) 是实数，\(m \in \mathbf{N}\) 是满足 \(m < k - n/2\) 的整数。空间 \(\mathrm{S}^k(\mathbf{R}^n)\) 含于 \(\mathcal{C}^m(\mathbf{R}^n)\) （利用柯西–施瓦茨不等式，验证 \(x \mapsto x^\alpha \mathscr{F} f(x)\) 属于 \(\mathrm{L}^1(\mathbf{R}^n)\) 对所有 \(\alpha \in \mathbf{N}^n\) （满足 \(|\alpha| \leqslant m\) ）。然后利用傅里叶反演公式证明 \(\mathscr{F}(f)\) 在 \(\mathbf{R}\) 上可微。）

\(g)\) 设 \(k > n/2\) 是实数，\(m \in \mathbf{N}\) 是满足 \(m < k - n/2\) 的整数。\(\mathrm{S}_{\mathrm{U}}^{k}(\mathbf{R}^{n})\) 的元素到 U 的限制属于 \(\mathcal{C}^{m}(\mathrm{U})\) （对每个 \(x \in \mathrm{U}\) ，验证得到一个良定义的函数 \(g \in \mathcal{C}^{m}(\mathrm{U})\) ，它由 \(g(x) = (\varphi_{x}f)(x)\) 定义，其中 \(\varphi_{x} \in \mathcal{D}(\mathbf{R}^{n})\) 是支撑在 U 中且在 \(x\) 的一个邻域中等于 1 的函数；最后验证 \(f\) 到 U 的限制等于 \(g\) 。）

\begin{enumerate}
\def\labelenumi{\alph{enumi})}
\setcounter{enumi}{7}
\tightlist
\item
  设 \(g \in \mathcal{D}'(\mathbf{R}^n)\) 。若 \(f \in \mathcal{D}'(\mathbf{R}^n)\) 满足 \(\Delta(f) = g\) 且 \(g\) 到 U 的限制属于 \(\mathcal{C}^\infty(\mathrm{U})\) ，则 \(f\) 到 U 的限制属于 \(\mathcal{C}^\infty(\mathrm{U})\) （“椭圆正则性定理”）。
\end{enumerate}

\begin{enumerate}
\def\labelenumi{\arabic{enumi})}
\setcounter{enumi}{27}
\tightlist
\item
  设 \(n\) 与 \(m\) 是整数，\(\mathrm{U} \subset \mathbf{R}^n\) 与 \(\mathrm{V} \subset \mathbf{R}^m\) 是开集。\(a)\) 存在唯一的单射线性映射 \(u\) ，它从 \(\mathcal{D}(\mathrm{U}) \otimes \mathcal{D}(\mathrm{V})\) 到 \(\mathcal{D}(\mathrm{U} \times \mathrm{V})\) 中，使得 \(u(\varphi \otimes \psi)\) 是函数 \((x, y) \mapsto \varphi(x)\psi(y)\) 对所有 \((\varphi, \psi) \in \mathcal{D}(\mathrm{U}) \times \mathcal{D}(\mathrm{V})\) 。此外，\(u\) 的像在 \(\mathcal{D}(\mathrm{U} \times \mathrm{V})\) 中稠密。我们把空间 \(\mathcal{D}(\mathrm{U}) \otimes \mathcal{D}(\mathrm{V})\) 等同于 \(\mathcal{D}(\mathrm{U} \times \mathrm{V})\) 的一个子空间，借助映射 \(u\) 。
\end{enumerate}

\begin{enumerate}
\def\labelenumi{\alph{enumi})}
\setcounter{enumi}{1}
\item
  对 U 上的所有广义函数 \(f\) 与 V 上的所有广义函数 \(g\) ，存在 U × V 上唯一的广义函数 \(f \otimes g\) ，使得 \(\langle f \otimes g, \varphi \otimes \psi \rangle = \langle f, \varphi \rangle \langle g, \psi \rangle\) 对所有 \((\varphi, \psi) \in \mathcal{D}(\mathrm{U}) \times \mathcal{D}(\mathrm{V})\) 成立。
\item
  设 \(f \in \mathcal{D}(\mathrm{U})\) 且 \(g \in \mathcal{D}(\mathrm{V})\) 。对每个 \(\eta \in \mathcal{D}(\mathrm{U} \times \mathrm{V})\) 以及每个 \(x \in U\)，用 \(\eta_x\) 表示映射 \(y \mapsto \eta(x, y)\)。我们有 \(\eta_x \in \mathcal{D}(\mathrm{V})\) ，并且映射 \(\varphi\) （它把每个 \(x \in U\) 映为 \(\langle f, \eta_x \rangle\) ）属于 \(\mathcal{D}(\mathrm{U})\) 且满足
\end{enumerate}

\[
\langle f \otimes g, \eta \rangle = \langle f, \varphi \rangle .
\]

\begin{enumerate}
\def\labelenumi{\arabic{enumi})}
\setcounter{enumi}{28}
\tightlist
\item
  设 \(n\) 与 \(m\) 是整数，\(\mathrm{U} \subset \mathbf{R}^n\) 与 \(\mathrm{V} \subset \mathbf{R}^m\) 是开集。我们沿用上一个习题的记号。设 \(u\) 是从 \(\mathcal{D}(\mathrm{U})\) 到 \(\mathcal{D}'(\mathrm{V})\) 中的连续线性映射。
\end{enumerate}

\begin{enumerate}
\def\labelenumi{\alph{enumi})}
\item
  存在唯一的连续线性形式 \(\widetilde{\kappa}\) ，它定义在子空间 \(\mathcal{D}(\mathrm{U})\otimes \mathcal{D}(\mathrm{V})\) （\(\mathcal{D}(\mathrm{U}\times \mathrm{V})\) 的）上，使得 \(\widetilde{\kappa} (\varphi \otimes \psi) = \langle u(\varphi),\psi \rangle\) 对所有 \((\varphi ,\psi)\) （\(\mathcal{D}(\mathrm{U})\times \mathcal{D}(\mathrm{V})\) 中的）成立。
\item
  存在广义函数 \(\kappa(u)\) （在 \(U \times V\) 上）延拓 \(\widetilde{\kappa}\) ，并且映射 \(u \mapsto \kappa(u)\) （从 \(\mathcal{L}(\mathcal{D}(U), \mathcal{D}'(V))\) 到 \(\mathcal{D}'(U \times V)\) 中的）是连续的。
\item
  映射 \(\kappa\) 是双射（“施瓦茨核定理”；为证明 \(\kappa\) 的满射性，给定 U × V 上的广义函数 \(\eta\) ，把 \(u(\varphi)\) 定义为 V 上的广义函数 \(\psi \mapsto \langle \eta, \varphi \otimes \psi \rangle\) 。）
\end{enumerate}

\begin{enumerate}
\def\labelenumi{\arabic{enumi})}
\setcounter{enumi}{29}
\tightlist
\item
  设 \(n \in \mathbf{N}\) 。我们用 \((x_1, \ldots, x_n)\) 表示 \(\mathbf{R}^n\) 的典范基的对偶基。
\end{enumerate}

\begin{enumerate}
\def\labelenumi{\alph{enumi})}
\tightlist
\item
  设 \(c \geqslant 0\) 与 r \textgreater{} 0 是实数。存在在 \(R^{n}\) 中 0 的开邻域 U 上定义且解析的函数 u 使得
\end{enumerate}

\[
\left(r - \sum_ {j = 1} ^ {n - 1} x _ {j} - u\right) \partial_ {n} u - c r \sum_ {j = 1} ^ {n - 1} \partial_ {j} u = c r
\]

在 U 上成立，并且 \(u(x,0)=0\) 对 \(x\in\mathbf{R}^{n-1}\) （满足 \((x,0)\in\mathrm{U}\) ）。（先考虑 n=2 的情形。）

\begin{enumerate}
\def\labelenumi{\alph{enumi})}
\setcounter{enumi}{1}
\tightlist
\item
  设 U 是 \(\mathbf{R}^{n}\) 中 0 的开邻域。设 \((a_{1},\ldots,a_{n-1},b)\) 是从 \(\mathrm{U}\times\mathbf{R}\) 到 \(\mathbf{C}\) 的一组函数，它们在 \((0,0)\in\mathrm{U}\times\mathbf{R}\) 的一个邻域中解析。我们用 Q 表示从 U 上解析函数所成的向量空间到其自身的映射，使得
\end{enumerate}

\[
\mathrm{Q} (u) = \partial_ {n} u - \sum_ {j = 1} ^ {n - 1} a _ {j} (x _ {1}, \dots , x _ {n - 1}, u) \partial_ {j} u + b (x _ {1}, \dots , x _ {n - 1}, u).
\]

至多存在一个函数 \(v\) ，它在 0 的开邻域 \(V\) （含于 \(U\) ）中定义且解析，使得 \(Q(v) = 0\) 并且 \(v(x,0) = 0\) 对所有 \(x \in \mathbf{R}^{n-1}\) （满足 \((x,0) \in V\) ）。（验证这样的解 \(v\) 在 0 处的泰勒展开的系数被唯一确定。）

\begin{enumerate}
\def\labelenumi{\alph{enumi})}
\setcounter{enumi}{2}
\tightlist
\item
  设 \(f\) 与 \(g\) 是 \(\mathbf{C}[[z_1, \ldots, z_n]]\) 中的形式幂级数。称 \(f\) 优控 \(g\) ，并记作 \(f \ll g\) ，如果
\end{enumerate}

\[
f = \sum_ {\alpha \in \mathbf {N} ^ {n}} a _ {\alpha} z ^ {\alpha}, \qquad g = \sum_ {\alpha \in \mathbf {N} ^ {n}} b _ {\alpha} z ^ {\alpha}
\]

其中 \(|a_{\alpha}| \leqslant b_{\alpha}\) 对所有 \(\alpha \in \mathbf{N}^{n}\) 。若形式级数 g 在 \(R^{n}\) 中 0 的一个邻域中收敛，则 f 亦然。

\(d)\) 若存在在 0 的邻域中解析的一组函数 \((\widetilde{a}_1, \ldots, \widetilde{a}_{n-1}, \widetilde{b})\) ，使得 \(a_j \ll \widetilde{a}_j\) 对 \(1 \leqslant j \leqslant n-1\) 且 \(b \ll \widetilde{b}\) （即在 0 处的相应泰勒展开满足这些条件的意义下），并且若方程

\[
\partial_ {n} v - \sum_ {j = 1} ^ {n - 1} \widetilde {a} _ {j} (x _ {1}, \dots , x _ {n - 1}, v) \partial_ {j} v + \widetilde {b} (x _ {1}, \dots , x _ {n - 1}, v) = 0
\]

关于未知函数 \(v\) 在 0 的一个邻域中有解析解，则方程 \(\mathrm{Q}(v) = 0\) 在 0 的一个邻域中有唯一的解析解。

\begin{enumerate}
\def\labelenumi{\alph{enumi})}
\setcounter{enumi}{4}
\tightlist
\item
  方程 \(Q(v) = 0\) 存在唯一的在 0 的邻域中解析的解（“柯西–柯瓦列夫斯卡娅定理”）。
\end{enumerate}

\begin{enumerate}
\def\labelenumi{\arabic{enumi})}
\setcounter{enumi}{30}
\tightlist
\item
  我们用 L 表示微分算子
\end{enumerate}

\[
\mathrm{L} = - i \partial_ {1} + \partial_ {2} - 2 (x _ {1} + i x _ {2}) \partial_ {3}
\]

它是 \(\mathbf{R}^{3}\) 上阶为 1 的算子，其中 \((x_{1}, x_{2}, x_{3})\) 是 \(\mathbf{R}^{3}\) 的典范基的对偶基，诸 \(x_{i}\) 被视为 \(\mathbf{R}^{3}\) 上的多项式（“卢伊算子”）。

\begin{enumerate}
\def\labelenumi{\alph{enumi})}
\item
  设 U 是 \(\mathbf{R}^3\) 中 0 的开邻域。存在函数 \(\varphi \in \mathcal{D}(\mathrm{U})\) 使得 \(\varphi(0) = 1\) 并且 \(\mathrm{L}(\varphi)\) 在含于 U 的 0 的邻域 V 中为零。（利用柯西–柯瓦列夫斯卡娅定理。）
\item
  设 \(u \in \mathcal{C}^{\infty}(\mathrm{U})\) 是满足 \(\mathrm{L}(u) = 0\) 且 \(u(0) = 0\) 的函数。假设存在 \(\delta > 0\) 使得 \(\mathcal{R}(u(x)) < -2\delta\) 对 \(x \in U - V\) 成立。我们令 \(v_{t} = \varphi \exp(t(u + \delta))\)。我们有 \(v_{t} \in \mathcal{C}^{\infty}(\mathrm{U})\) ，并且 \(\mathrm{L}(v_{t})\) 在 \(\mathcal{D}(\mathrm{U})\) 中趋于 0 当 \(t \to +\infty\) 时。
\item
  设 F 是由这样的函数 \(f \in \mathcal{C}^{\infty}(U)\) 组成的空间：\(\partial^{\alpha}f\) 对每个 \(\alpha \in N^{3}\) 在 U 上有界。空间 F 赋有半范数 \(q_{\alpha}(f) = \sup_{x \in U} |\partial^{\alpha}f|\) 后是弗雷歇空间。
\end{enumerate}

\(d)\) 设 \(\psi\in\mathcal{D}(\mathbf{R}^{3})\) 。用 \(f_{t}\) 表示函数 \(x\mapsto t^{3}\psi(tx)\) 到 U 的限制。我们有 \(f_{t}\in\mathrm{F}\) 对 \(t\geqslant1\) ，并且存在实数 \(c\geqslant0\) ，使得 \(q_{\alpha}(f_{t})\leqslant ct^{|\alpha|+3}\) 对所有 \(\alpha\in\mathbf{N}^{3}\) 以及每个 \(t\geqslant1\) 。

\begin{enumerate}
\def\labelenumi{\alph{enumi})}
\setcounter{enumi}{4}
\tightlist
\item
  对每个 \(t \geqslant 1\) ，我们用 \(\lambda_{t}\) 表示 F 上的线性形式 \(f \mapsto \int_{U} fv_{t}\)。它是连续的，并且对充分大的 t，我们有
\end{enumerate}

\[
\lambda_ {t} (f _ {t}) = e ^ {\delta t} \left(\int_ {\mathrm{U}} \psi (x) \exp \left(\sum_ {j = 1} ^ {3} \xi_ {j} x _ {j}\right) d x _ {1} d x _ {2} d x _ {3} + \varepsilon (t)\right)
\]

其中 \(\xi_{j} = \partial_{j}u(0)\) ，并且 \(\varepsilon (t)\to 0\) 当 \(t\rightarrow +\infty\)。

\(f)\) 假设 \((\xi_j) \neq 0\) ；存在 \(f \in \mathrm{F}\) 使得积分 \(\int_{\mathrm{U}} f v_t\) 当 \(t \to +\infty\) 时不是有界的。（若该论断不真，则利用上一问题与巴拿赫–斯坦豪斯定理得出矛盾。）

\(g)\) 诸函数 \(v_{t}\) 在 \(\mathcal{D}'(\mathrm{U})\) 中不收敛于 0 当 \(t\to+\infty\) 时。\(h)\) 设 \(\xi=(0,0,\xi_{3})\in\mathbf{R}^{3}\) 满足 \(\xi_{3}<0\)。存在 \(\mathbf{R}^{3}\) 中 0 的邻域 U 以及函数 \(u\in\mathcal{C}^{\infty}(\mathrm{U})\) ，使得 \(\mathrm{L}(u)=0\) 且 \(u(0)=0\) ，并满足问题 \(b)\) 与 \(f)\) 的条件。（利用柯西–柯瓦列夫斯卡娅定理。）

\begin{enumerate}
\def\labelenumi{\roman{enumi})}
\tightlist
\item
  设 \((v_{t})\) 是 \(\mathcal{D}(\mathrm{U})\) 中满足 \(\mathrm{L}(v_{t})\to0\) 的函数，并设 \(f\in\mathcal{D}(\mathrm{U})\) 使得 \(\int_{\mathrm{U}}fv_{t}\) 不是有界的当 \(t\to+\infty\) 时（问题 g)）；方程 Lv=f 没有解 \(v\in\mathcal{D}'(\mathrm{U})\) （“卢伊定理”）。（注意 L 的形式伴随与 L 一致。）
\end{enumerate}

\begin{enumerate}
\def\labelenumi{\arabic{enumi})}
\setcounter{enumi}{31}
\tightlist
\item
  设 \(n \geqslant 1\) 是整数。我们用 \(\Delta\) 表示 \(\mathcal{D}'(\mathbf{R}^n)\) 上的拉普拉斯算子，它定义为
\end{enumerate}

\[
\Delta (f) = - \sum_ {i = 1} ^ {n} \partial_ {i} ^ {2} f.
\]

设 u 是 \(R^{n}\) 上的局部可积函数，使得 \(\Delta(u)\) 是局部可积函数。

设 \(v\in \mathcal{L}^{\infty}(\mathbf{R}^{n})\) 是由下式定义的函数：

\[
v (x) = \left\{ \begin{array}{l l} 0 & \text { if } u (x) = 0, \\ \overline {{u (x)}} / | u (x) | & \text { if } u (x) \neq 0. \end{array} \right.
\]

\begin{enumerate}
\def\labelenumi{\alph{enumi})}
\item
  对 \(\varepsilon > 0\)，令 \(u_{\varepsilon} = (|u|^2 + \varepsilon^2)^{1/2}\) 。证明广义函数 \(-\Delta(u_{\varepsilon}) - \mathcal{R}(\overline{u}/u_{\varepsilon} \Delta(u))\) 是正的。（先考虑 \(u \in \mathcal{C}^{\infty}(\mathbf{R}^{n})\) 的情形，再用卷积处理一般情形。）
\item
  广义函数 \(-\Delta(|u|) - \mathcal{R}(v\Delta(u))\) 是正的（“加藤不等式”）。
\end{enumerate}

\begin{enumerate}
\def\labelenumi{\arabic{enumi})}
\setcounter{enumi}{32}
\tightlist
\item
  设 \(m \geqslant 1\) 是整数。
\end{enumerate}

\begin{enumerate}
\def\labelenumi{\alph{enumi})}
\tightlist
\item
  存在唯一的线性映射 \(A^{+}\) （相应地 \(A^{-}\) ），它从 \(\mathrm{H}^{m}(\mathbf{R})\) 到 \(\mathrm{H}^{m-1}(\mathbf{R})\) ，使得
\end{enumerate}

\[
\mathrm{A} ^ {+} f (x) = f ^ {\prime} (x) - x f (x) \quad (\text { resp.   A } ^ {-} f (x) = f ^ {\prime} (x) + x f (x))
\]

对所有 \(f \in \mathcal{D}(\mathbf{R})\) 以及每个 \(x \in \mathbf{R}\) 成立。

\begin{enumerate}
\def\labelenumi{\alph{enumi})}
\setcounter{enumi}{1}
\tightlist
\item
  映射 \(A^{+}\) 是指标为 1 的弗雷德霍姆算子，而 \(A^{-}\) 是指标为 -1 的弗雷德霍姆算子。（考虑埃尔米特函数在 \(A^{+}\) 与 \(A^{-}\) 下的像，参见 II, p. 263 的习题 3。）
\end{enumerate}
% semantic-fragments: legacy-input-seam
\relax{}
